\chapter{Cronograma de ejecucion}
\thispagestyle{memoria}\pagestyle{memoria}

\section{Plazo referencial}
Para un grifo de esta magnitud se estima un plazo referencial de ejecucion de
8 semanas, con la siguiente distribucion de hitos. Es un orden academico
referencial; el cronograma definitivo depende de la aprobacion del proyecto y
de los permisos sectoriales.

\begin{center}
\begin{tabularx}{\linewidth}{@{}L l@{}}
\toprule
Hito & Semana \\
\midrule
Obra civil y canalizaciones enterradas & 1--2 \\
Pozos de puesta a tierra y equipotencialidad & 2 \\
Alimentadores, subalimentadores y tableros & 3--4 \\
Circuitos derivados (alumbrado, tomacorrientes, fuerza) & 4--5 \\
Equipos de playa (surtidores, STP, compresor, bombas) & 5--6 \\
Grupo electrogeno, ATS, UPS y pararrayo & 6--7 \\
Pruebas, verificaciones y documentacion & 7--8 \\
\bottomrule
\end{tabularx}
\end{center}

\section{Hitos de control}
Los hitos criticos son: energizacion de tableros, puesta en servicio del
grupo electrogeno, prueba del paro de emergencia y medicion de la resistencia
de puesta a tierra. Cada hito requiere registro fotografico y verificacion
por el residente y el supervisor segun corresponda.
